\chapter{Discussion}

\section{Overall interpretation}
\label{sec:ch8-overall}

The results support a bounded account of efficient closed-set VQA on this GQA
development partition. Compact trainable heads can use frozen representations
to combine strong question-side predictive information with a real visual
contribution. At 40k, the question-only head reaches 0.45797 and fusion reaches
0.53840. The intervention evidence then shows that the multimodal heads depend
on receiving the correct image; it does not establish visual grounding or
reasoning. This use of frozen features with a small closed-answer classifier
has precedent on a different VQA setting \citep{deuser2022less}, but the
dataset and protocol differences preclude a numerical comparison.

Across the implemented systems, representation choice, interaction design and
answer-space coverage all materially affect the observed results. Additional
downstream capacity alone does not consistently explain the strongest
configurations. This is not a global ranking of those design choices: the
top-100 and top-1000 evaluations define different supported-answer tasks, and
cross-family comparisons often change several factors together. The following
sections therefore interpret the controlled contrasts first and use the less
closely matched comparisons only as qualified context.

\section{Representations, interaction and latent reasoning}
\label{sec:ch8-representation}

The RQ1 evidence gives explicit interaction a clear but limited role. Fusion
exceeds concatenation by +0.01442 at 40k, yet the capacity controls show that
part of this difference follows from the wider head. At the larger matched
budget, the fusion-minus-concat-wide effect is +0.00578 with interval
$[-0.00116,0.01226]$, so that residual contrast is not resolved clearly.
This separation of interaction from an additive approximation is conceptually
related to the question posed by EMAP \citep{hessel2020emap}; here it is tested
constructively with capacity controls. The later narrowing of the
fusion--concatenation gap remains descriptive because no clustered interval is
available for that scaling family.

Representation choice also matters, but the encoder comparison is not
factorial. At 250k, SigLIP fusion at 0.59388 is descriptively competitive with
the latent reasoner at 0.59580 in the tested configurations.
Encoder and interface widths, architecture and other configuration details
co-vary, and the seed structures differ. The observation is therefore neither
a paired causal comparison nor evidence that a stronger representation can
generally replace downstream reasoning capacity.

The reasoner's overall accuracy rises across the tested scales, and its
same-seed advantage over fusion is positive in every stored seed at the two
larger scales. This remains configuration-level evidence because architecture,
token access and capacity change together. Learned latent bottlenecks have
architectural precedents in Perceiver and BLIP-2
\citep{jaegle2021perceiver,li2023blip2}, but those lineages do not identify the
source of this project's accuracy pattern. More importantly, the fixed-prior
four-or-more-step deficit persists and its paired difference from fusion is
unresolved at every tested scale. The implemented changes did not eliminate
that profile, but they do not localise its cause.

The question-only result provides an internal check on what this statistic
means. At the largest scale it has a smaller pooled deficit than the
multimodal systems despite substantially lower overall accuracy. A lower
deficit therefore describes a less uneven fixed-prior step profile, not a
better model or stronger reasoning ability. Conversely, the reasoner's higher
overall accuracy at the larger scales does not establish that its persistent
multi-step difficulty has been removed.

That restraint is consistent with the wider evidence base. Binding information
can remain linearly accessible in unimodal CLIP features
\citep{koishigarina2026bow}, while work with much larger decoders locates
feature-use limitations differently or finds compensation for degraded visual
context \citep{fu2025hidden,takishita2025compensate}. These are empirical
counterpoints, not matched comparisons: together they caution against calling
compositional information absent or assigning this deficit to one component.
The GQA program length is only a proxy for reasoning depth
\citep{hudson2019gqa}, and the buckets have different answer priors, a known
concern for GQA evaluation \citep{kervadec2021gqaood}.

The token diagnostics narrow the interpretation further. The latent queries
place approximately 0.26--0.38 attention mass on CLS across seeds, compared
with a uniform reference near 0.02, and removing CLS reduces accuracy. These
observations show disproportionate attention and functional sensitivity to the
pooled token. They do not show that most attention is pooled, that the model
explicitly re-reads CLS, that attention weights trace causal information flow,
or how and why the token is used internally.

\section{What SLM placement suggests}
\label{sec:ch8-placement}

The question-side comparisons must be read together. A1 exceeds its
architecture-matched random control A1r with positive directional evidence at
both tested scales, so frozen SmolLM2 pretraining contributes within that
interface. At the same time, A1 remains below the CLIP-token A0p reference at
both scales. The latter contrast is secondary and representation-confounded:
it changes the representation and interface together and cannot establish
intrinsic CLIP superiority.

The answer-side evidence has a different pattern. The 135M
pretrained-minus-random contrast is negative directional at the lower scale and uncertain or
mixed at the larger scale. Both 360M contrasts and the second-order scale
interaction are also uncertain or mixed. Thus no reliable positive
answer-side pretraining advantage was detected under these conditions, while
the zero-containing intervals establish neither equivalence nor absence of an
effect.

Several explanations are plausible. The projected representation spaces may
differ in compatibility; contextual question structure is available before
multimodal interaction only on the question side; and a frozen language model
placed after the reasoner may make limited use of its pretraining under a
closed candidate-likelihood readout and restricted adaptation. Minimal bridges
into frozen language models and comparisons among connector families provide
architectural context for these possibilities
\citep{merullo2023limber,vallaeys2024depalm}, but they do not explain this
result. The experiments do not distinguish among the hypotheses.

The resulting claim is therefore interface-specific rather than universal
under the implemented interfaces and tested configurations. It is a bounded
cross-experiment synthesis of three separately run experiments, whose family
and checkpoint-selection contexts differ, not a factorial causal placement
study. It neither establishes a generally optimal placement nor rules out a
positive answer-side contribution in another regime.

\section{Accuracy--efficiency trade-offs}
\label{sec:ch8-efficiency}

Chapter 7 adds a measurement boundary to the accuracy results. Under the E7b
warm-serial protocol, concat, fusion and the top-1000 product head form the
canonical frontier on common-denominator raw-distribution accuracy and latency.
The higher-accuracy token-level configurations are not automatically preferred
on those two axes. Some additional components coincide with higher measured
latency, but not all accuracy gains arise from added compute: the representation
swap and answer-space experiment also move results without following a single
capacity ordering. Nor does this pattern imply that more computation causes
more accuracy. Trainable and total loaded parameter counts likewise provide
capacity context rather than latency proxies.

The E9 comparison remains contextual. Its stored statement makes the
SmolVLM-500M constrained row approximately 24.7 times slower in warm serial
latency than the re-timed head on the same node, but the comparison is
latency-only.
The E9 artefact does not restate the head checkpoint identity, so its head
accuracy and timing remain unpaired. E7b and E9 are not merged because the
bridge gate failed, and no matched accuracy--latency leaderboard follows.
Energy, power, carbon and monetary cost were not measured, so no claim is made
about those quantities.

\section{Limitations}
\label{sec:ch8-limitations}

The evidential boundary is development-only. Checkpoint and recipe choices use
development evidence, and no held-out result is reported; the F2 evaluation
remains unauthorised and pending. The task is one partition of GQA under closed
answer classification. Top-100 and top-1000 systems differ in answer support,
training rows and in-vocabulary denominator, so their supported-view
accuracies cannot be treated as one metric. These constraints limit
generalisation to other datasets, answer spaces and generative settings.

Several comparisons also remain configuration-level. Reasoner and fusion
differ simultaneously in architecture, token access and trainable capacity;
the SigLIP comparison changes encoder and width; and the SLM evidence covers a
small set of frozen family sizes, two placements, two training scales and
differing selection rules across experiment families. The placement synthesis
therefore does not isolate a universal pretraining effect, and the descriptive
SigLIP--reasoner proximity does not isolate representation from downstream
capacity. These limitations narrow the claims without undoing the controlled
within-family contrasts.

The measurement and compositional analyses have separate limits. E7b and E9
were run on different nodes and remain separate; E9 head pairing is unresolved,
and comparable serial latency is unavailable for the answer-side R1 readout
and the E10 pair. The step analysis uses program length as an in-distribution
proxy, combines unequal buckets under fixed priors and is sensitive to the
question/answer distribution. It describes a persistent relative difficulty
profile, not observed internal reasoning steps or a localised mechanism. The
experiments therefore leave several causes open rather than identifying a
single bottleneck; proposals for resolving them belong to Chapter 9.

\section{Overall implications}
\label{sec:ch8-implications}

Within the implemented systems and governed development protocol, frozen
representations with compact heads provide useful closed-set VQA performance
and a favourable measured latency region. Explicit interaction helps most
clearly in the low-data comparison, representation changes can be competitive
with much larger downstream configurations, and latent capacity improves
overall accuracy at larger scales without eliminating the multi-step deficit.
Frozen language-model pretraining contributes differently across the tested
interfaces rather than simply through its presence.

Taken together, the evidence favours disciplined selection among
representation, interaction, answer-space and interface choices, not a single
ranked design rule. The simplest global heads occupy the measured frontier,
while the more elaborate systems clarify what extra capacity and token access
do not by themselves establish. Whether these development findings generalise
is unresolved; evaluation and design proposals are deferred to Chapter 9.
